
%(BEGIN_QUESTION)
% Copyright 2006, Tony R. Kuphaldt, released under the Creative Commons Attribution License (v 1.0)
% This means you may do almost anything with this work of mine, so long as you give me proper credit

{\it Nernst-ligningen} er hovedformelen for å beregne spenning over {\it enhver} ioneselektiv membran, inkludert pH-elektroder. Generell form:

$$V = {{R T} \over {nF}} \ln \left({C_1 \over C_2}\right)$$

\noindent
Der,

$V$ = Spenning over membranen, volt (V)

$R$ = Universell gasskonstant (8.315 J/mol$\cdot$K)

$T$ = Absolutt temperatur, Kelvin (K)

$n$ = Antall elektroner per utvekslet ion

$F$ = Faraday-konstant (96,485 C/mol e$^{-}$)

$C_1$ = Ionkonsentrasjon i måleløsning ($M$)

$C_2$ = Ionkonsentrasjon i referanseløsning ($M$)

\vskip 10pt

Nernst-ligningen kan også skrives med 10-logaritmer:

$$V = {{2.303 R T} \over {nF}} \log \left({C_1 \over C_2}\right)$$

Beregn pH for følgende konsentrasjoner ved 25$^{o}$ C, og bruk Nernst-ligningen til å beregne spenning over glassmembranen i en pH-elektrode. Husk at intern KCl-referanse er buffret til [H$^{+}$] = 1.00 $\times$ $10^{-7}$ $M$. Bruk $n=1$ for hydrogenioner:

% No blank lines allowed between lines of an \halign structure!
% I use comments (%) instead, so that TeX doesn't choke.

$$\vbox{\offinterlineskip
\halign{\strut
\vrule \quad\hfil # \ \hfil & 
\vrule \quad\hfil # \ \hfil & 
\vrule \quad\hfil # \ \hfil \vrule \cr
\noalign{\hrule}
%
% First row
[H$^{+}$] i løsning & pH & V \cr
%
\noalign{\hrule}
%
% Another row
5.83 $\times$ $10^{-5}$ $M$ &  &  \cr
%
\noalign{\hrule}
%
% Another row
1 $\times$ $10^{-5}$ $M$ &  &  \cr
%
\noalign{\hrule}
%
% Another row
1 $\times$ $10^{-6}$ $M$ &  &  \cr
%
\noalign{\hrule}
%
% Another row
1 $\times$ $10^{-7}$ $M$ &  &  \cr
%
\noalign{\hrule}
%
% Another row
1 $\times$ $10^{-8}$ $M$ &  &  \cr
%
\noalign{\hrule}
%
% Another row
7.02 $\times$ $10^{-9}$ $M$ &  &  \cr
%
\noalign{\hrule}
%
% Another row
1 $\times$ $10^{-9}$ $M$ &  &  \cr
%
\noalign{\hrule}
%
% Another row
1 $\times$ $10^{-10}$ $M$ &  &  \cr
%
\noalign{\hrule}
} % End of \halign 
}$$ % End of \vbox

Etter noen beregninger vil du se et mønster. Finn proporsjonaliteten mellom generert spenning og pH, og uttrykk den enkelt.

\underbar{file i00618}
%(END_QUESTION)





%(BEGIN_ANSWER)

% No blank lines allowed between lines of an \halign structure!
% I use comments (%) instead, so that TeX doesn't choke.

$$\vbox{\offinterlineskip
\halign{\strut
\vrule \quad\hfil # \ \hfil & 
\vrule \quad\hfil # \ \hfil & 
\vrule \quad\hfil # \ \hfil \vrule \cr
\noalign{\hrule}
%
% First row
[H$^{+}$] i løsning & pH & V \cr
%
\noalign{\hrule}
%
% Another row
5.83 $\times$ $10^{-5}$ $M$ & 4.23 & 163.6 mV \cr
%
\noalign{\hrule}
%
% Another row
1 $\times$ $10^{-5}$ $M$ & 5 & 118.3 mV \cr
%
\noalign{\hrule}
%
% Another row
1 $\times$ $10^{-6}$ $M$ & 6 & 59.16 mV \cr
%
\noalign{\hrule}
%
% Another row
1 $\times$ $10^{-7}$ $M$ & 7 & 0 mV \cr
%
\noalign{\hrule}
%
% Another row
1 $\times$ $10^{-8}$ $M$ & 8 & $-59.16$ mV \cr
%
\noalign{\hrule}
%
% Another row
7.02 $\times$ $10^{-9}$ $M$ & 8.15 & $-68.25$ mV \cr
%
\noalign{\hrule}
%
% Another row
1 $\times$ $10^{-9}$ $M$ & 9 & $-118.3$ mV \cr
%
\noalign{\hrule}
%
% Another row
1 $\times$ $10^{-10}$ $M$ & 10 & $-177.5$ mV \cr
%
\noalign{\hrule}
} % End of \halign 
}$$ % End of \vbox

Forskjellen i pH mellom måleløsning og intern referanse uttrykkes som en enkel DC-spenning.

\vskip 10pt

Helning = 59.17 mV per pH-enhet avvik fra 7.0 pH

\vskip 10pt

Temperaturkompensering kreves ved høynøyaktig pH-måling for å kompensere temperaturleddet i Nernst-ligningen. Dette gjelder membranen, ikke løsningen. RTD i sonden korrigerer altså pH/spenningsforholdet i elektroden, men kan ikke kompensere for reelle ionaktivitetsendringer i løsningen. For det må man kjenne løsningens unike temperatur/[H$^{+}$]-sammenheng.

\vskip 10pt

Den mer korrekte Nernst-formen er:

$$V = {{R T} \over {nF}} \ln \left({a_1 \over a_2}\right)$$

\noindent
Der,

$a_1$ = Ioneaktivitet i måleløsning

$a_2$ = Ioneaktivitet i referanseløsning

\vskip 10pt

Uten kjemiske reaksjoner som ``binder opp'' ioner, er aktivitet ($a$) og konsentrasjon ($C$) proporsjonale ($a = \gamma C$). Aktivitetskoeffisienten ($\gamma$) kan likevel variere og gjøre uttrykket mer komplekst.

%(END_ANSWER)





%(BEGIN_NOTES)


%INDEX% Kjemi, elektro-: Nernst-ligningen
%INDEX% Kjemi, pH: molaritetsberegning
%INDEX% Måling, analytisk: pH

%(END_NOTES)

